% from main.tex
%%%%%%%%%%  scan idx 249  |  folio 236  %%%%%%%%%%
\textbf{Corollaire 1.12.} Dans la situation de (1.11), si \(g\) est pseudo-cohérent, on a, pour \(r\in\mathbb Z\cup\{-\infty\}\),
\[
  D^{-}(f)_{r\text{-}\coh}=D^{-}(h)_{r\text{-}\coh};
\]
en particulier (faisant \(X=Y\)), les notions de \(r\)-pseudo-cohérence sur \(Y\), « au sens absolu » ou rel. à \(g\), coïncident.

\textbf{Corollaire 1.13.} Dans la situation de (1.11), si \(f\) est pseudo-cohérent, le foncteur
\[
  \mathcal O_X\otimes_Y^{\mathbf L}.=Lf^*:D^{-}(Y)\longrightarrow D^{-}(X)
\]
induit, pour \(r\in\mathbb Z\cup\{-\infty\}\), un foncteur
\[
  Lf^*:D^{-}(g)_{r\text{-}\coh}\longrightarrow D^{-}(h)_{r\text{-}\coh}.
\]

\textbf{Corollaire 1.14.} Dans la situation de (1.11), si \(f\) est \(r\)-pseudo-cohérent et \(g\) est \(s\)-pseudo-cohérent (\(r,s\in\mathbb Z\cup\{-\infty\}\)), alors \(h\) est \(\operatorname{Sup}(r,s)\)-pseudo-cohérent. Si \(h\) est \(r\)-pseudo-cohérent et \(g\) est \((r+1)\)-pseudo-cohérent, \(f\) est \(r\)-pseudo-cohérent.

\textbf{Remarque 1.14.1.} En particulier, pour \(r=s=-1\), on retrouve (cf. (1.2.1 c))) les assertions (ii) et (v) de (EGA IV 1.4.3).

\section*{2. Le théorème de finitude}

\textbf{Conjecture 2.1.} Soient \(S\) un schéma, \(f:X\to Y\) un \(S\)-morphisme propre de schémas localement de type fini sur \(S\), \(E\in\Ob D^+(X)\). Alors si \(E\) est pseudo-cohérent rel. à \(S\) (1.2), il en est de même de \(\mathbf Rf_*(E)\).

\textbf{Théorème 2.2.} La conjecture précédente est vraie dans les cas suivants :
\[
\begin{gathered}
  \text{(i) } S\text{ est localement noethérien};\\
  \text{(ii) } f\text{ est projectif}.
\end{gathered}
\]

